\section{Conclusion}
\label{sec:conclusion}

This thesis formalized carbon-aware data-center scheduling as a deterministic linear program that minimizes operational carbon by jointly choosing when and where flexible load runs across five structurally diverse grid regions. Driven by hourly carbon-intensity and carbon-free-energy signals over a two-year horizon, and subject to seven operational constraints, the model was validated through a behavioural perturbation test and an integration check on real data, and then evaluated in a rolling-window backtest. Under unconstrained flexibility it reduced carbon by up to 78\,\% relative to a no-optimization baseline, and under realistic operational constraints it saved 31\,\% while remaining the only method in the comparison that both attained near-ceiling savings and provably respected every constraint.

Beyond the headline saving, two findings carry the main message of the work. First, the sensitivity analysis identified geographic flexibility as the dominant determinant of achievable savings, exceeding every other parameter by an order of magnitude, which indicates that the binding constraint on carbon-aware scheduling is how much load may move across regions rather than how it is shaped in time. Second, the spatio-temporal decomposition showed that routing load to cleaner regions accounts for almost all of the saving, whereas shifting load in time alone, with load confined to a high-carbon home region, can increase carbon. Together these results reframe the practical question for operators from when to run flexible work to where to run it, and they show that the value of an exact optimum lies less in beating a well-tuned heuristic on a given dataset than in providing a provably optimal reference and a guarantee that operational limits are respected.

Several limitations qualify these conclusions and motivate further work. The model uses average carbon intensity rather than locational marginal emissions, which attribute carbon to the generator that responds to an incremental load and can diverge from average intensity in large grids \citep{cote2024locational, lindberg2022geographic}; adopting a marginal signal is a natural extension. The negative value of time-shifting alone is specific to PJM as the home region, and the cross-region timing comparison rests on a five-point diagnostic rather than a statistical test, so both should be read as indicative rather than conclusive. The formulation is also deterministic and assumes that carbon signals are known over the planning window, whereas a deployment must act on forecasts; a stochastic or distributionally robust extension would quantify the carbon penalty of forecast error \citep{li2024uncertainty, hall2024probabilistic}. Finally, the model optimizes aggregate load against operational carbon alone, leaving task-level scheduling, embodied hardware emissions, and water consumption outside its scope \citep{gupta2021chasing, li2023thirsty}.

These limitations suggest a clear agenda for future work, including a marginal-emissions objective, a robust or stochastic treatment of forecast uncertainty, a multi-objective formulation that weighs carbon against water and other footprints, and an extension to precedence-constrained, task-level workloads \citep{lechowicz2025pcaps, attenni2024shifting}. Even within its present scope, however, the model offers a validated and provably optimal account of carbon-aware spatio-temporal scheduling, and a clear empirical answer to which form of flexibility actually drives carbon savings.
